% Problem record op_1ab8b10386bddd66. This TeX file is derived from
% database/problems_json/op_1ab8b10386bddd66.json by scripts/sync-tex.mjs; edit the JSON record
% and rerun the script rather than editing this file.
\section{Visible quantum compression with free shared randomness at the Holevo rate}
\label{sec:op_1ab8b10386bddd66}

\paragraph{Problem.}
Does free shared classical randomness make the Holevo information achievable for visible compression of every finite mixed-state ensemble? Let $p$ be a probability distribution with full support on a finite alphabet $\mathcal X$, and let $W_x$ be density operators on a finite-dimensional system $B$. Alice knows the independent and identically distributed label sequence $x^n$, with probability $p^n(x^n)$; the target is $W_{x^n}:=\bigotimes_{t=1}^nW_{x_t}$. Alice and Bob share a finite random variable $S_n$, independent of $x^n$, with arbitrary distribution $q_n$ and no bound on its size. For each seed $s$, Alice prepares a density operator $\tau_{x^n,s}$ on a quantum message system $M_n$ and sends it through a noiseless quantum channel. Bob applies a completely positive trace-preserving map $\mathcal D_{n,s}:\mathcal L(M_n)\to\mathcal L(B^{\otimes n})$. No initially shared entanglement or additional communication is available. The reproduced state and block error are
\begin{equation}
 \widehat W_{x^n}:=\sum_s q_n(s)\mathcal D_{n,s}(\tau_{x^n,s}),\qquad
 \varepsilon_n:=\frac12\sum_{x^n}p^n(x^n)\|W_{x^n}-\widehat W_{x^n}\|_1.
 \label{eq:h15-error}
\end{equation}
Thus the shared seed is averaged before taking the trace norm in Eq.~\eqref{eq:h15-error}. Define the rate, in qubits per signal, and the ensemble Holevo information by
\begin{equation}
 R_{\mathrm{vis},q,\mathrm{cr}}(p,W):=\inf_{\{\text{codes}:\varepsilon_n\to0\}}\limsup_{n\to\infty}\frac1n\log_2\dim M_n,
 \qquad
 \chi(p,W):=S\!\left(\sum_xp_xW_x\right)-\sum_xp_xS(W_x),
 \label{eq:h15-rate}
\end{equation}
where $S$ is the von Neumann entropy in bits. Is $R_{\mathrm{vis},q,\mathrm{cr}}(p,W)=\chi(p,W)$ in Eq.~\eqref{eq:h15-rate} for every ensemble, or is there a finite ensemble with a strict gap?

\subsection*{Status}

Unsolved

\subsection*{Source}

This is the free-randomness communication-rate part of Winter's Conjecture~23, p.~11, with the visible quantum model defined in Eqs.~(15)--(16) of \sourcecite{ref:h15-winter}{Win02}. The original conjecture additionally specifies a shared-randomness rate of $\sum_xp_xS(W_x)$; the present question leaves that resource unrestricted and is consequently weaker.

\subsection*{Progress}

\begin{itemize}
  \item The converse $R_{\mathrm{vis},q,\mathrm{cr}}(p,W)\ge\chi(p,W)$ follows from Holevo information and data processing, consistent with Winter's quantum lower bounds \sourcecite{ref:h15-winter}{Win02}, Theorems~20--21. For the shared-randomness model specifically, independence of $S_n$ and the labels gives $I(X^n;M_nS_n)=I(X^n;M_n\mid S_n)\le\log_2\dim M_n$; vanishing block error and entropy continuity then yield the claimed asymptotic converse.

  \item For the analogous protocol with a classical message, Nator and Pereg's Corollary~3 gives
\begin{equation}
 R_{\mathrm{vis},c,\mathrm{cr}}(p,W)=\min_{P(u\mid x),\theta_u:\,W_x=\sum_uP(u\mid x)\theta_u}I(X;U),
 \label{eq:h15-classical}
\end{equation}
where $\theta_u$ are density operators, $P_{XU}(x,u)=p_xP(u\mid x)$, and one may take $|\mathcal U|\le |\mathcal X|^2(\dim B)^2+1$. Their joint-state trace-norm criterion, Eqs.~(3)--(4), equals Eq.~\eqref{eq:h15-error} by the trace norm of block-diagonal operators \sourcecite{ref:h15-nator}{NP24}. A classical message can be encoded in orthogonal quantum states, so Eq.~\eqref{eq:h15-classical} is an upper bound on the quantum rate. For commuting $W_x$, this value is $\chi(p,W)$, settling that subclass.

  \item Without shared randomness, Hayashi's main theorem, Eq.~(8), identifies the visible qubit rate with $E_P^\infty(\omega_{XB})$, where $\omega_{XB}=\sum_xp_x|x\rangle\langle x|\otimes W_x$, $E_P$ is the minimum entropy across a bipartition of a purification, and $E_P^\infty(\omega):=\lim_{n\to\infty}E_P(\omega^{\otimes n})/n$. Ignoring shared randomness therefore gives the upper bound $R_{\mathrm{vis},q,\mathrm{cr}}\le E_P^\infty(\omega_{XB})$. For pure signal states, Eq.~(10) gives $E_P^\infty(\omega_{XB})=S(\sum_xp_xW_x)=\chi(p,W)$, so that subclass is also settled \sourcecite{ref:h15-hayashi}{Hay06}.
\end{itemize}

\subsection*{References}

\begin{enumerate}[label={},leftmargin=4.8em,labelsep=0.5em]
  \footnotesize
  \item[\textup{[Win02]}]\label{ref:h15-winter}
  A. Winter, ``Compression of Sources of Probability Distributions and Density Operators,'' arXiv preprint (2002). \href{https://arxiv.org/abs/quant-ph/0208131}{arXiv:quant-ph/0208131}.

  \item[\textup{[NP24]}]\label{ref:h15-nator}
  H. Nator and U. Pereg, ``Coordination Capacity for Classical-Quantum Correlations,'' arXiv preprint (2024). \href{https://arxiv.org/abs/2404.18297}{arXiv:2404.18297}.

  \item[\textup{[Hay06]}]\label{ref:h15-hayashi}
  M. Hayashi, ``Optimal Visible Compression Rate for Mixed States Is Determined by Entanglement of Purification,'' \emph{Physical Review A} \textbf{73}, 060301(R) (2006). \href{https://doi.org/10.1103/PhysRevA.73.060301}{doi:10.1103/PhysRevA.73.060301}; \href{https://arxiv.org/abs/quant-ph/0511267}{arXiv:quant-ph/0511267}.
\end{enumerate}

\subsection*{Comment}

The remaining issue is quantum achievability at the Holevo rate for general noncommuting mixed signals. The classical communication formula in Eq.~\eqref{eq:h15-classical} settles a different resource restriction and does not supply this quantum coding theorem. Likewise, entanglement-assisted reverse Shannon protocols use a resource excluded here. This question is distinct from additivity of entanglement of purification: that problem controls the unassisted visible rate, whereas the present task supplies shared randomness \sourcecite{ref:h15-hayashi}{Hay06}. The placement of the average in Eq.~\eqref{eq:h15-error} is essential: averaging the trace-norm errors of individual seeds would allow a good seed to be fixed separately at each blocklength, removing any benefit from shared randomness. Literature audited on 9 September 2026. The cited primary definitions and theorem passages were checked, and no full resolution was identified in the later sources examined.

\subsection*{Field}

Quantum Communication

\subsection*{Topic}

Quantum source coding; Channel simulation

\subsection*{ID}

\texttt{op\_1ab8b10386bddd66}
